\section{Alcance, riesgos y criterios de salida}

Dentro del alcance: AMM constante, precio CEX externo, roles representativos aprobados por el autor
con políticas simplificadas, mapa discreto local trazable, E0--E3, reproducibilidad, experimentos
sintéticos y comparación limitada con hechos de mercado.

Fuera del alcance: predicción rentable, exchanges operativos, réplicas LOB completas,
consenso blockchain, MEV avanzado, lending, bridges, aprendizaje por refuerzo y
garantías generales sobre todos los mercados. El MVP actual no establece validez
empírica del mapa. En su primera comprobación reservada, la etiqueta de Schur coincidió
con el resultado mayoritario de contracción en solo 3 de 8 grupos. Riesgo principal:
que el mapa reducido no aproxime adecuadamente al ABM fuera de una región local.
Mitigación: revisar las reglas de agentes, separar selección y confirmación, registrar
configuraciones, contrastar ablaciones y reportar desacuerdos como resultado. Otros
riesgos: sobreparametrización, dependencia de semillas, liquidez simplificada, confusión
entre estabilidad local y global y sesgo en calibración.

\section{Discusión y límites de inferencia}

El análisis debe distinguir tres niveles: propiedades mecánicas deducidas del AMM ideal,
diagnóstico local del mapa reducido y patrones de un ABM discreto estocástico sujetos a
validación. El mapa actual clasificó correctamente solo 3 de 8 resultados mayoritarios
reservados bajo un cambio acotado de semilla y magnitud del shock; su condición de Schur
no garantiza estabilidad en el ABM ni en mercados reales. La ejecución de selección
tampoco sustituye una evaluación con datos públicos. El trabajo futuro deberá informar
los regímenes donde la aproximación reduzca el error y aquellos donde falle por
saturación, límites de capital, shocks persistentes o cambios de liquidez.

\section{Conclusiones y trabajo futuro}

El proyecto cuenta ahora con un ejecutor determinista de mecanismos, un modelo
estocástico de selección emparejada y un diagnóstico reservado del mapa reducido. Las
comprobaciones deterministas muestran que la capacidad de arbitraje cambia la velocidad
de corrección y que menos reservas aumentan el impacto de una venta fija en el CPMM
configurado. La comparación estocástica muestra que el efecto del LP adaptativo depende
de la dirección del shock. La primera prueba H1c coincide con solo 3 de 8 grupos de
resultados mayoritarios, por lo que el mapa actual no es un clasificador fiable. La
siguiente fase consiste en cerrar los roles de agentes, ampliar las regiones operativas,
prerregistrar criterios confirmatorios y evaluar un conjunto público acotado. Ningún
resultado actual demuestra utilidad en un mercado real.
